% --- Archon physics formalization source begin ---
% archon:physics
% archon:covers PhyXMiniProblems/problem_phyx_mini_0395.lean
% archon:source-report reports/phyx_mini/problem_phyx_mini_0395.source.json

\chapter{Physics problem phyx\_mini\_0395}
\label{ch:PhyXMiniProblems_problem_phyx_mini_0395}

\paragraph{Problem source.}
\#\# Physical scenario

Two cylinders are connected by a piston. Cylinder A is used as a hydraulic lift and pumped up to 500 kPa. The piston mass is 25 kg, and there is standard gravity.

\#\# Question

What is the gas pressure in cylinder B?

\#\# Answer choices

- A: A: 4MPa

- B: B: 6MPa

- C: C: 154kPa

- D: D: 10.2kPa

\#\# Figure caption (auxiliary; use the image as primary evidence)

The image is a diagram depicting a hydraulic system with two pistons labeled "A" and "B." These pistons are enclosed in a vertical, cylindrical configuration. 

- Piston A is at the bottom and is larger, with a diameter (\textbackslash{}(D\_A\textbackslash{})) of 100 mm.
- Piston B is at the top, with a smaller diameter (\textbackslash{}(D\_B\textbackslash{})) of 25 mm.
- Both pistons are shown within shaded cylindrical sections, which likely represent the chambers in which these pistons move.

The system is pressurized, indicated by the notation \textbackslash{}(P\_0 = 100 \textbackslash{}, \textbackslash{}text\{kPa\}\textbackslash{}) next to piston B.

There is an arrow labeled \textbackslash{}(g\textbackslash{}) pointing downward, likely representing gravitational force.

A pump symbol is connected to the lower chamber, suggesting the presence of fluid flow or pressure management within the system.

The relationships between the parts suggest a basic hydraulic mechanism where pressure applied on piston A can be transferred to piston B, or vice versa, facilitated by the pump.

\#\# Dataset metadata

Laws of Thermodynamics; Physical Model Grounding Reasoning, Multi-Formula Reasoning

\paragraph{Recorded answer/context.}
B: B: 6MPa

\paragraph{Figure/image path.}
/root/proposal\_for\_physic/science-mango/phyx\_mini\_run/phyx\_data/test\_image/395.png

\paragraph{Formalization target.}
create a compiling Lean file with sorry bodies at `PhyXMiniProblems/problem\_phyx\_mini\_0395.lean`. The Lean declarations must preserve the physical quantities, dimensions or dimensional roles, figure labels, governing-law hypotheses, and final relation expressed by this problem.
Use Mathlib/Physlib names found through LeanExplore where available. If a physics API is missing, introduce faithful local abstractions rather than scalar placeholder aliases.

\begin{theorem}[Physics formalization target]
\label{thm:physics:phyx_mini_0395:target}
\lean{PhyXMiniProblems.ProblemPhyXMini0395.gasPressureInCylinderB}
\uses{def:physics:phyx-mini-0395:phyxminiproblems-problemphyxmini0395-pressureinpascals, def:physics:phyx-mini-0395:phyxminiproblems-problemphyxmini0395-compoundpistonsetup, def:physics:phyx-mini-0395:phyxminiproblems-problemphyxmini0395-matchesprimaryfigure, def:physics:phyx-mini-0395:phyxminiproblems-problemphyxmini0395-matchesproblemdata, def:physics:phyx-mini-0395:phyxminiproblems-problemphyxmini0395-hasphysicalparameters, def:physics:phyx-mini-0395:phyxminiproblems-problemphyxmini0395-satisfiescircularfacegeometry, def:physics:phyx-mini-0395:phyxminiproblems-problemphyxmini0395-matchesidealstaticmodel, def:physics:phyx-mini-0395:phyxminiproblems-problemphyxmini0395-satisfiesstaticforcebalance, def:physics:phyx-mini-0395:phyxminiproblems-problemphyxmini0395-isuniqueclosestanswer}
The assigned autoformalize agent should translate this physics problem into Lean declarations in the covered file, with theorem and lemma proof bodies written as `by sorry`.
\end{theorem}
\begin{proof}
This is an autoformalization task, not a proof task. Produce statements that can later be proved by the physics prover without weakening the physical model.
\end{proof}

% --- Archon declaration topology begin ---
\section{Declaration topology}
\begin{definition}[Length Quantity]
  \label{def:physics:phyx-mini-0395:phyxminiproblems-problemphyxmini0395-lengthquantity}
  \lean{PhyXMiniProblems.ProblemPhyXMini0395.LengthQuantity}
  A nonnegative physical length, independent of the unit used to read it
\end{definition}
\begin{proof}
  This is the declared model definition of Length Quantity.
\end{proof}
\begin{definition}[Mass Quantity]
  \label{def:physics:phyx-mini-0395:phyxminiproblems-problemphyxmini0395-massquantity}
  \lean{PhyXMiniProblems.ProblemPhyXMini0395.MassQuantity}
  A nonnegative physical mass
\end{definition}
\begin{proof}
  This is the declared model definition of Mass Quantity.
\end{proof}
\begin{definition}[Acceleration Quantity]
  \label{def:physics:phyx-mini-0395:phyxminiproblems-problemphyxmini0395-accelerationquantity}
  \lean{PhyXMiniProblems.ProblemPhyXMini0395.AccelerationQuantity}
  A nonnegative physical acceleration magnitude
\end{definition}
\begin{proof}
  This is the declared model definition of Acceleration Quantity.
\end{proof}
\begin{definition}[Area Quantity]
  \label{def:physics:phyx-mini-0395:phyxminiproblems-problemphyxmini0395-areaquantity}
  \lean{PhyXMiniProblems.ProblemPhyXMini0395.AreaQuantity}
  A nonnegative physical cross-sectional area
\end{definition}
\begin{proof}
  This is the declared model definition of Area Quantity.
\end{proof}
\begin{definition}[Pressure Quantity]
  \label{def:physics:phyx-mini-0395:phyxminiproblems-problemphyxmini0395-pressurequantity}
  \lean{PhyXMiniProblems.ProblemPhyXMini0395.PressureQuantity}
  A physical absolute pressure
\end{definition}
\begin{proof}
  This is the declared model definition of Pressure Quantity.
\end{proof}
\begin{definition}[length In Meters]
  \label{def:physics:phyx-mini-0395:phyxminiproblems-problemphyxmini0395-lengthinmeters}
  \lean{PhyXMiniProblems.ProblemPhyXMini0395.lengthInMeters}
  \uses{def:physics:phyx-mini-0395:phyxminiproblems-problemphyxmini0395-lengthquantity}
  Coherent-SI metre readout of a physical length
\end{definition}
\begin{proof}
  This is the declared model definition of length In Meters.
\end{proof}
\begin{definition}[length In Millimeters]
  \label{def:physics:phyx-mini-0395:phyxminiproblems-problemphyxmini0395-lengthinmillimeters}
  \lean{PhyXMiniProblems.ProblemPhyXMini0395.lengthInMillimeters}
  \uses{def:physics:phyx-mini-0395:phyxminiproblems-problemphyxmini0395-lengthquantity, def:physics:phyx-mini-0395:phyxminiproblems-problemphyxmini0395-lengthinmeters}
  Millimetre readout used by the two diameter labels in the figure
\end{definition}
\begin{proof}
  This is the declared model definition of length In Millimeters.
\end{proof}
\begin{definition}[mass In Kilograms]
  \label{def:physics:phyx-mini-0395:phyxminiproblems-problemphyxmini0395-massinkilograms}
  \lean{PhyXMiniProblems.ProblemPhyXMini0395.massInKilograms}
  \uses{def:physics:phyx-mini-0395:phyxminiproblems-problemphyxmini0395-massquantity}
  Coherent-SI kilogram readout of a physical mass
\end{definition}
\begin{proof}
  This is the declared model definition of mass In Kilograms.
\end{proof}
\begin{definition}[acceleration In Meters Per Second Squared]
  \label{def:physics:phyx-mini-0395:phyxminiproblems-problemphyxmini0395-accelerationinmeterspersecondsquared}
  \lean{PhyXMiniProblems.ProblemPhyXMini0395.accelerationInMetersPerSecondSquared}
  \uses{def:physics:phyx-mini-0395:phyxminiproblems-problemphyxmini0395-accelerationquantity}
  Coherent-SI metre-per-second-squared readout of an acceleration
\end{definition}
\begin{proof}
  This is the declared model definition of acceleration In Meters Per Second Squared.
\end{proof}
\begin{definition}[area In Square Meters]
  \label{def:physics:phyx-mini-0395:phyxminiproblems-problemphyxmini0395-areainsquaremeters}
  \lean{PhyXMiniProblems.ProblemPhyXMini0395.areaInSquareMeters}
  \uses{def:physics:phyx-mini-0395:phyxminiproblems-problemphyxmini0395-areaquantity}
  Coherent-SI square-metre readout of a physical area
\end{definition}
\begin{proof}
  This is the declared model definition of area In Square Meters.
\end{proof}
\begin{definition}[pressure In Pascals]
  \label{def:physics:phyx-mini-0395:phyxminiproblems-problemphyxmini0395-pressureinpascals}
  \lean{PhyXMiniProblems.ProblemPhyXMini0395.pressureInPascals}
  \uses{def:physics:phyx-mini-0395:phyxminiproblems-problemphyxmini0395-pressurequantity}
  Coherent-SI pascal readout of an absolute pressure
\end{definition}
\begin{proof}
  This is the declared model definition of pressure In Pascals.
\end{proof}
\begin{definition}[pressure In Kilopascals]
  \label{def:physics:phyx-mini-0395:phyxminiproblems-problemphyxmini0395-pressureinkilopascals}
  \lean{PhyXMiniProblems.ProblemPhyXMini0395.pressureInKilopascals}
  \uses{def:physics:phyx-mini-0395:phyxminiproblems-problemphyxmini0395-pressurequantity, def:physics:phyx-mini-0395:phyxminiproblems-problemphyxmini0395-pressureinpascals}
  Kilopascal readout used by the data printed in the problem
\end{definition}
\begin{proof}
  This is the declared model definition of pressure In Kilopascals.
\end{proof}
\begin{definition}[Cylinder Label]
  \label{def:physics:phyx-mini-0395:phyxminiproblems-problemphyxmini0395-cylinderlabel}
  \lean{PhyXMiniProblems.ProblemPhyXMini0395.CylinderLabel}
  The two cylinders labelled in the supplied image
\end{definition}
\begin{proof}
  This is the declared model definition of Cylinder Label.
\end{proof}
\begin{definition}[Chamber Contents]
  \label{def:physics:phyx-mini-0395:phyxminiproblems-problemphyxmini0395-chambercontents}
  \lean{PhyXMiniProblems.ProblemPhyXMini0395.ChamberContents}
  Physical contents associated with a labelled cylinder
\end{definition}
\begin{proof}
  This is the declared model definition of Chamber Contents.
\end{proof}
\begin{definition}[Piston Surface]
  \label{def:physics:phyx-mini-0395:phyxminiproblems-problemphyxmini0395-pistonsurface}
  \lean{PhyXMiniProblems.ProblemPhyXMini0395.PistonSurface}
  Exposed surfaces of the stepped piston that carry vertical pressure loads
\end{definition}
\begin{proof}
  This is the declared model definition of Piston Surface.
\end{proof}
\begin{definition}[Vertical Direction]
  \label{def:physics:phyx-mini-0395:phyxminiproblems-problemphyxmini0395-verticaldirection}
  \lean{PhyXMiniProblems.ProblemPhyXMini0395.VerticalDirection}
  Vertical direction of an arrow or force component in the image
\end{definition}
\begin{proof}
  This is the declared model definition of Vertical Direction.
\end{proof}
\begin{definition}[Primary Hydraulic Figure]
  \label{def:physics:phyx-mini-0395:phyxminiproblems-problemphyxmini0395-primaryhydraulicfigure}
  \lean{PhyXMiniProblems.ProblemPhyXMini0395.PrimaryHydraulicFigure}
  \uses{def:physics:phyx-mini-0395:phyxminiproblems-problemphyxmini0395-cylinderlabel, def:physics:phyx-mini-0395:phyxminiproblems-problemphyxmini0395-pistonsurface, def:physics:phyx-mini-0395:phyxminiproblems-problemphyxmini0395-verticaldirection}
  Qualitative labels and placements read directly from image 395.png
\end{definition}
\begin{proof}
  This is the declared model definition of Primary Hydraulic Figure.
\end{proof}
\begin{definition}[Compound Piston Setup]
  \label{def:physics:phyx-mini-0395:phyxminiproblems-problemphyxmini0395-compoundpistonsetup}
  \lean{PhyXMiniProblems.ProblemPhyXMini0395.CompoundPistonSetup}
  \uses{def:physics:phyx-mini-0395:phyxminiproblems-problemphyxmini0395-lengthquantity, def:physics:phyx-mini-0395:phyxminiproblems-problemphyxmini0395-massquantity, def:physics:phyx-mini-0395:phyxminiproblems-problemphyxmini0395-accelerationquantity, def:physics:phyx-mini-0395:phyxminiproblems-problemphyxmini0395-areaquantity, def:physics:phyx-mini-0395:phyxminiproblems-problemphyxmini0395-pressurequantity, def:physics:phyx-mini-0395:phyxminiproblems-problemphyxmini0395-cylinderlabel, def:physics:phyx-mini-0395:phyxminiproblems-problemphyxmini0395-chambercontents, def:physics:phyx-mini-0395:phyxminiproblems-problemphyxmini0395-primaryhydraulicfigure}
  The stepped piston and its independent physical observables. In particular, the pressure in cylinder B is an unconstrained physical pressure here; its requested numerical value is not built into this structure
\end{definition}
\begin{proof}
  This is the declared model definition of Compound Piston Setup.
\end{proof}
\begin{definition}[Matches Primary Figure]
  \label{def:physics:phyx-mini-0395:phyxminiproblems-problemphyxmini0395-matchesprimaryfigure}
  \lean{PhyXMiniProblems.ProblemPhyXMini0395.MatchesPrimaryFigure}
  \uses{def:physics:phyx-mini-0395:phyxminiproblems-problemphyxmini0395-compoundpistonsetup}
  The labels and spatial roles visible in the supplied primary figure
\end{definition}
\begin{proof}
  This is the declared model definition of Matches Primary Figure.
\end{proof}
\begin{definition}[Matches Problem Data]
  \label{def:physics:phyx-mini-0395:phyxminiproblems-problemphyxmini0395-matchesproblemdata}
  \lean{PhyXMiniProblems.ProblemPhyXMini0395.MatchesProblemData}
  \uses{def:physics:phyx-mini-0395:phyxminiproblems-problemphyxmini0395-lengthinmillimeters, def:physics:phyx-mini-0395:phyxminiproblems-problemphyxmini0395-massinkilograms, def:physics:phyx-mini-0395:phyxminiproblems-problemphyxmini0395-accelerationinmeterspersecondsquared, def:physics:phyx-mini-0395:phyxminiproblems-problemphyxmini0395-pressureinkilopascals, def:physics:phyx-mini-0395:phyxminiproblems-problemphyxmini0395-compoundpistonsetup}
  Numerical information supplied by the prose and image. The value 196133 / 20000 m/s² is standard gravity 9.80665 m/s². The pressure in cylinder B is deliberately absent
\end{definition}
\begin{proof}
  This is the declared model definition of Matches Problem Data.
\end{proof}
\begin{definition}[Has Physical Parameters]
  \label{def:physics:phyx-mini-0395:phyxminiproblems-problemphyxmini0395-hasphysicalparameters}
  \lean{PhyXMiniProblems.ProblemPhyXMini0395.HasPhysicalParameters}
  \uses{def:physics:phyx-mini-0395:phyxminiproblems-problemphyxmini0395-lengthinmeters, def:physics:phyx-mini-0395:phyxminiproblems-problemphyxmini0395-massinkilograms, def:physics:phyx-mini-0395:phyxminiproblems-problemphyxmini0395-accelerationinmeterspersecondsquared, def:physics:phyx-mini-0395:phyxminiproblems-problemphyxmini0395-areainsquaremeters, def:physics:phyx-mini-0395:phyxminiproblems-problemphyxmini0395-pressureinpascals, def:physics:phyx-mini-0395:phyxminiproblems-problemphyxmini0395-compoundpistonsetup}
  Positivity and ordering assumptions selecting the depicted physical branch
\end{definition}
\begin{proof}
  This is the declared model definition of Has Physical Parameters.
\end{proof}
\begin{definition}[Satisfies Circular Face Geometry]
  \label{def:physics:phyx-mini-0395:phyxminiproblems-problemphyxmini0395-satisfiescircularfacegeometry}
  \lean{PhyXMiniProblems.ProblemPhyXMini0395.SatisfiesCircularFaceGeometry}
  \uses{def:physics:phyx-mini-0395:phyxminiproblems-problemphyxmini0395-lengthinmeters, def:physics:phyx-mini-0395:phyxminiproblems-problemphyxmini0395-areainsquaremeters, def:physics:phyx-mini-0395:phyxminiproblems-problemphyxmini0395-cylinderlabel, def:physics:phyx-mini-0395:phyxminiproblems-problemphyxmini0395-compoundpistonsetup}
  Both load-bearing piston faces are circular. This is a generic geometry law relating each physical area to its corresponding figure diameter; it contains no pressure value and no answer choice
\end{definition}
\begin{proof}
  This is the declared model definition of Satisfies Circular Face Geometry.
\end{proof}
\begin{definition}[Matches Ideal Static Model]
  \label{def:physics:phyx-mini-0395:phyxminiproblems-problemphyxmini0395-matchesidealstaticmodel}
  \lean{PhyXMiniProblems.ProblemPhyXMini0395.MatchesIdealStaticModel}
  \uses{def:physics:phyx-mini-0395:phyxminiproblems-problemphyxmini0395-compoundpistonsetup}
  Qualitative idealizations used in the static pressure-force model
\end{definition}
\begin{proof}
  This is the declared model definition of Matches Ideal Static Model.
\end{proof}
\begin{definition}[upward Hydraulic Force In Newtons]
  \label{def:physics:phyx-mini-0395:phyxminiproblems-problemphyxmini0395-upwardhydraulicforceinnewtons}
  \lean{PhyXMiniProblems.ProblemPhyXMini0395.upwardHydraulicForceInNewtons}
  \uses{def:physics:phyx-mini-0395:phyxminiproblems-problemphyxmini0395-areainsquaremeters, def:physics:phyx-mini-0395:phyxminiproblems-problemphyxmini0395-pressureinpascals, def:physics:phyx-mini-0395:phyxminiproblems-problemphyxmini0395-compoundpistonsetup}
  Upward hydraulic pressure force on the large lower face, in newtons
\end{definition}
\begin{proof}
  This is the declared model definition of upward Hydraulic Force In Newtons.
\end{proof}
\begin{definition}[downward Gas Force In Newtons]
  \label{def:physics:phyx-mini-0395:phyxminiproblems-problemphyxmini0395-downwardgasforceinnewtons}
  \lean{PhyXMiniProblems.ProblemPhyXMini0395.downwardGasForceInNewtons}
  \uses{def:physics:phyx-mini-0395:phyxminiproblems-problemphyxmini0395-areainsquaremeters, def:physics:phyx-mini-0395:phyxminiproblems-problemphyxmini0395-pressureinpascals, def:physics:phyx-mini-0395:phyxminiproblems-problemphyxmini0395-compoundpistonsetup}
  Downward gas pressure force on the small upper face, in newtons
\end{definition}
\begin{proof}
  This is the declared model definition of downward Gas Force In Newtons.
\end{proof}
\begin{definition}[annular Shoulder Area In Square Meters]
  \label{def:physics:phyx-mini-0395:phyxminiproblems-problemphyxmini0395-annularshoulderareainsquaremeters}
  \lean{PhyXMiniProblems.ProblemPhyXMini0395.annularShoulderAreaInSquareMeters}
  \uses{def:physics:phyx-mini-0395:phyxminiproblems-problemphyxmini0395-areainsquaremeters, def:physics:phyx-mini-0395:phyxminiproblems-problemphyxmini0395-compoundpistonsetup}
  Area of the exposed annular shoulder, read in square metres
\end{definition}
\begin{proof}
  This is the declared model definition of annular Shoulder Area In Square Meters.
\end{proof}
\begin{definition}[downward Atmospheric Force In Newtons]
  \label{def:physics:phyx-mini-0395:phyxminiproblems-problemphyxmini0395-downwardatmosphericforceinnewtons}
  \lean{PhyXMiniProblems.ProblemPhyXMini0395.downwardAtmosphericForceInNewtons}
  \uses{def:physics:phyx-mini-0395:phyxminiproblems-problemphyxmini0395-pressureinpascals, def:physics:phyx-mini-0395:phyxminiproblems-problemphyxmini0395-compoundpistonsetup, def:physics:phyx-mini-0395:phyxminiproblems-problemphyxmini0395-annularshoulderareainsquaremeters}
  Downward atmospheric force on the annular shoulder, in newtons
\end{definition}
\begin{proof}
  This is the declared model definition of downward Atmospheric Force In Newtons.
\end{proof}
\begin{definition}[piston Weight In Newtons]
  \label{def:physics:phyx-mini-0395:phyxminiproblems-problemphyxmini0395-pistonweightinnewtons}
  \lean{PhyXMiniProblems.ProblemPhyXMini0395.pistonWeightInNewtons}
  \uses{def:physics:phyx-mini-0395:phyxminiproblems-problemphyxmini0395-massinkilograms, def:physics:phyx-mini-0395:phyxminiproblems-problemphyxmini0395-accelerationinmeterspersecondsquared, def:physics:phyx-mini-0395:phyxminiproblems-problemphyxmini0395-compoundpistonsetup}
  Downward gravitational force on the piston, in newtons
\end{definition}
\begin{proof}
  This is the declared model definition of piston Weight In Newtons.
\end{proof}
\begin{definition}[Satisfies Static Force Balance]
  \label{def:physics:phyx-mini-0395:phyxminiproblems-problemphyxmini0395-satisfiesstaticforcebalance}
  \lean{PhyXMiniProblems.ProblemPhyXMini0395.SatisfiesStaticForceBalance}
  \uses{def:physics:phyx-mini-0395:phyxminiproblems-problemphyxmini0395-compoundpistonsetup, def:physics:phyx-mini-0395:phyxminiproblems-problemphyxmini0395-upwardhydraulicforceinnewtons, def:physics:phyx-mini-0395:phyxminiproblems-problemphyxmini0395-downwardgasforceinnewtons, def:physics:phyx-mini-0395:phyxminiproblems-problemphyxmini0395-downwardatmosphericforceinnewtons, def:physics:phyx-mini-0395:phyxminiproblems-problemphyxmini0395-pistonweightinnewtons}
  The ideal vertical force balance for the compound piston. Cylinder A pushes upward over the large face. The gas in B, atmospheric pressure on the annular shoulder, and the piston's weight act downward. This governing law contains no numerical value for the unknown gas pressure
\end{definition}
\begin{proof}
  This is the declared model definition of Satisfies Static Force Balance.
\end{proof}
\begin{definition}[Answer Choice]
  \label{def:physics:phyx-mini-0395:phyxminiproblems-problemphyxmini0395-answerchoice}
  \lean{PhyXMiniProblems.ProblemPhyXMini0395.AnswerChoice}
  Labels of the four answers printed with the problem
\end{definition}
\begin{proof}
  This is the declared model definition of Answer Choice.
\end{proof}
\begin{definition}[displayed Pressure In Pascals]
  \label{def:physics:phyx-mini-0395:phyxminiproblems-problemphyxmini0395-displayedpressureinpascals}
  \lean{PhyXMiniProblems.ProblemPhyXMini0395.displayedPressureInPascals}
  \uses{def:physics:phyx-mini-0395:phyxminiproblems-problemphyxmini0395-answerchoice}
  Absolute pressure in pascals printed beside each answer label
\end{definition}
\begin{proof}
  This is the declared model definition of displayed Pressure In Pascals.
\end{proof}
\begin{definition}[recorded Answer Choice]
  \label{def:physics:phyx-mini-0395:phyxminiproblems-problemphyxmini0395-recordedanswerchoice}
  \lean{PhyXMiniProblems.ProblemPhyXMini0395.recordedAnswerChoice}
  \uses{def:physics:phyx-mini-0395:phyxminiproblems-problemphyxmini0395-answerchoice}
  The answer label recorded in the supplied dataset
\end{definition}
\begin{proof}
  This is the declared model definition of recorded Answer Choice.
\end{proof}
\begin{definition}[Is Unique Closest Answer]
  \label{def:physics:phyx-mini-0395:phyxminiproblems-problemphyxmini0395-isuniqueclosestanswer}
  \lean{PhyXMiniProblems.ProblemPhyXMini0395.IsUniqueClosestAnswer}
  \uses{def:physics:phyx-mini-0395:phyxminiproblems-problemphyxmini0395-answerchoice, def:physics:phyx-mini-0395:phyxminiproblems-problemphyxmini0395-displayedpressureinpascals}
  A displayed answer is the unique closest choice to a derived pressure
\end{definition}
\begin{proof}
  This is the declared model definition of Is Unique Closest Answer.
\end{proof}
% --- Archon declaration topology end ---
% --- Archon physics formalization source end ---
